\chapter{Analyse et conception}

\section{Introduction}
Ce chapitre traduit les objectifs en une solution concrète. Il présente
l'analyse des besoins, l'architecture générale de la plateforme, la conception de
la chaîne de traitement, le modèle dimensionnel de l'entrepôt, la spécification
des ratios et du score de performance, ainsi que la stratégie de prévision.

\section{Analyse des besoins}
\subsection{Besoins fonctionnels}
La solution doit permettre de :
\begin{itemize}[noitemsep]
  \item importer et fiabiliser les états financiers des quatre sociétés ;
  \item générer les clôtures trimestrielles manquantes de la \sfbt ;
  \item calculer les ratios financiers et le score global de performance ;
  \item comparer les sociétés entre elles (benchmark) ;
  \item stocker les résultats dans un entrepôt interrogeable ;
  \item prévoir les valeurs financières futures ;
  \item restituer l'information sous forme de tableaux de bord.
\end{itemize}

\subsection{Besoins non fonctionnels}
\begin{itemize}[noitemsep]
  \item \textbf{fiabilité} : les calculs doivent être exacts et vérifiés ;
  \item \textbf{reproductibilité} : régénération complète par une commande ;
  \item \textbf{traçabilité} : conservation des libellés d'origine et
        signalement des anomalies ;
  \item \textbf{modularité} : ajout aisé de nouvelles sociétés ou de nouveaux
        ratios ;
  \item \textbf{portabilité} : exécution multiplateforme (Windows, Linux).
\end{itemize}

\section{Architecture générale de la solution}
La plateforme suit une architecture décisionnelle en couches, illustrée par la
figure~\ref{fig:archi}. Les fichiers sources sont d'abord traités par la chaîne
\ac{ETL} (nettoyage, unification, augmentation), produisant un jeu de données
unifié. Celui-ci alimente le moteur de calcul des ratios et du \ac{SGPI}, puis
l'entrepôt de données. Enfin, le module de prévision et les tableaux de bord
exploitent ces données.

\begin{figure}[H]
\centering
\begin{tikzpicture}[
  node distance=0.7cm and 1.1cm,
  box/.style={rectangle,rounded corners,draw=darkblue,very thick,
    fill=lightblue,align=center,minimum height=1.05cm,minimum width=2.6cm,
    font=\footnotesize},
  arr/.style={-{Latex[length=2.5mm]},very thick,darkblue},
]
\node[box] (src) {Sources Excel\\(4 sociétés)};
\node[box,right=of src] (etl) {ETL\\nettoyage,\\unification};
\node[box,right=of etl] (aug) {Augmentation\\(trimestres)};
\node[box,below=of aug] (ratio) {Ratios + SGPI\\(46 indicateurs)};
\node[box,left=of ratio] (dw) {Data Warehouse\\(schéma étoile)};
\node[box,left=of dw] (ml) {Prévision ML\\(2 modèles)};
\node[box,below=of dw] (resti) {Restitution\\Power BI};
\draw[arr] (src) -- (etl);
\draw[arr] (etl) -- (aug);
\draw[arr] (aug) -- (ratio);
\draw[arr] (ratio) -- (dw);
\draw[arr] (dw) -- (ml);
\draw[arr] (dw) -- (resti);
\draw[arr] (ml) |- (resti);
\end{tikzpicture}
\caption{Architecture générale de la plateforme décisionnelle.}
\label{fig:archi}
\end{figure}

\section{Conception de la chaîne de traitement (ETL)}
\subsection{Unification des sources}
Les quatre fichiers présentant des formats différents, un \textbf{schéma cible
unique} a été défini, sous forme de table longue (\textit{tidy}) dont les
principales colonnes sont : société, état financier, indicateur (libellé propre,
libellé brut et clé normalisée), date, période, valeur, devise, échelle, type de
donnée (réel ou estimé) et source.

La difficulté majeure réside dans l'\textbf{hétérogénéité des libellés} de la
\sfbt : un même concept est écrit de nombreuses façons. La conception prévoit une
\textbf{clé de normalisation} (passage en minuscules, suppression des accents, de
la ponctuation et des espaces) qui regroupe automatiquement les variantes d'un
même indicateur, condition indispensable à la reconstitution des séries
temporelles.

\subsection{Classification des états et fiabilisation}
Pour la \sfbt, dont la feuille maîtresse empile tous les états sans étiquette, un
\textbf{découpage positionnel} reclasse chaque ligne dans son état (bilan actif,
bilan passif, compte de résultat, flux de trésorerie, soldes intermédiaires de
gestion). Des règles de fiabilisation traitent les anomalies : suppression des
doublons stricts, et correction des \textbf{fuites de colonnes comparatives}
(valeurs d'une année recopiées par erreur sur l'année suivante).

\subsection{Augmentation des données}
La \sfbt ne publiant que les clôtures de juin et décembre, la conception prévoit
la \textbf{génération des clôtures de mars et septembre}. La méthode distingue la
nature des grandeurs :
\begin{itemize}[noitemsep]
  \item les \textbf{grandeurs de stock} (postes de bilan) sont obtenues par
        \textbf{interpolation linéaire} entre deux clôtures connues ;
  \item les \textbf{grandeurs de flux} (compte de résultat, flux), cumulées sur
        l'exercice, sont obtenues par \textbf{répartition au prorata} du temps.
\end{itemize}
Les valeurs ainsi produites sont explicitement marquées comme estimées.

\section{Conception de l'entrepôt de données}
L'entrepôt adopte un \textbf{schéma en étoile} (figure~\ref{fig:etoile}) : une
table de faits centrale, entourée de dimensions conformes. La granularité de la
table de faits est : une ligne par société, date et indicateur. Les dimensions
décrivent les axes d'analyse.

\begin{figure}[H]
\centering
\begin{tikzpicture}[
  dim/.style={rectangle,rounded corners,draw=midblue,thick,fill=lightblue,
    align=center,minimum height=1cm,minimum width=2.8cm,font=\footnotesize},
  fact/.style={rectangle,draw=espritred,very thick,fill=espritred!10,
    align=center,minimum height=1.4cm,minimum width=3.2cm,font=\footnotesize\bfseries},
  link/.style={thick,midblue},
]
\node[fact] (f) {Fait\\Ratios / États};
\node[dim,above=1.4cm of f] (t) {Dim\_Temps};
\node[dim,below=1.4cm of f] (i) {Dim\_Indicateur\\/ Dim\_Ratio};
\node[dim,left=2cm of f] (s) {Dim\_Societe};
\node[dim,right=2cm of f] (r) {Dim\_Ratio};
\draw[link] (f) -- (t);
\draw[link] (f) -- (i);
\draw[link] (f) -- (s);
\draw[link] (f) -- (r);
\end{tikzpicture}
\caption{Modèle dimensionnel en étoile de l'entrepôt de données.}
\label{fig:etoile}
\end{figure}

Quatre dimensions sont définies : \texttt{Dim\_Societe} (société, rôle, secteur,
devise), \texttt{Dim\_Temps} (date, année, période), \texttt{Dim\_Indicateur}
(poste comptable, état, nature stock/flux) et \texttt{Dim\_Ratio} (code, libellé,
catégorie, benchmark). Des \textbf{vues SQL} prêtes à l'emploi facilitent
l'interrogation et l'usage par \textit{Power BI}.

\section{Spécification des ratios et du score de performance}
Conformément au référentiel adopté, \textbf{46 ratios} sont spécifiés, répartis
en sept catégories : gestion des liquidités, des actifs, des passifs, des
ressources humaines, des risques, rentabilité et cycle d'exploitation. Chaque
ratio est défini par sa formule, son unité et sa fourchette de référence.

Le \textbf{\ac{SGPI}} agrège ces ratios. Chaque ratio clé reçoit une note de 0 à
5 selon une grille de barème ; les notes sont moyennées par catégorie, puis
pondérées selon le tableau~\ref{tab:poids} pour obtenir un score sur 100.

\begin{table}[H]
\centering
\caption{Pondération des catégories dans le score global (SGPI).}
\label{tab:poids}
\begin{tabular}{lc}
\toprule
\textbf{Catégorie} & \textbf{Poids} \\
\midrule
Rentabilité & 25~\% \\
Structure financière & 20~\% \\
Liquidité & 15~\% \\
Gestion des actifs & 15~\% \\
Cycle d'exploitation & 10~\% \\
Gestion des risques & 10~\% \\
Productivité RH & 5~\% \\
\midrule
\textbf{Total} & \textbf{100~\%} \\
\bottomrule
\end{tabular}
\end{table}

\section{Conception du module de prévision}
La prévision cible les principales grandeurs financières de la \sfbt (chiffre
d'affaires, total des actifs, capitaux propres, résultat net). La conception
retient une \textbf{mise au format supervisé} de la série : la cible est le
\textbf{taux de croissance} (et non la valeur brute), choix essentiel pour rendre
la série stationnaire et permettre aux deux modèles d'extrapoler correctement.
Les variables explicatives combinent la tendance temporelle, les retards et la
saisonnalité. Deux modèles — \textbf{régression linéaire} et \textbf{forêt
aléatoire} — sont comparés par \textit{backtest} chronologique.

\section{Conclusion}
Ce chapitre a défini l'architecture de la solution et la conception détaillée de
chacune de ses composantes : chaîne \ac{ETL}, entrepôt en étoile, spécification
des ratios et du score, et stratégie de prévision. Le chapitre suivant présente
la réalisation effective et les résultats obtenus.
